\chapter{Tropical and idempotent algebra}
\index{tropical algebra}\index{idempotent algebra}\index{graded reduction}
\label{app:tropical-algebra}

This appendix fixes the algebra and signs underlying tropical charts,
tropical spectra, and calibrated polyhedra.
The two recurring conventions are
\[
                       v(f)=-\log|f|
\]
and the min-plus semiring.
Passing from \(u=-\log|f|\) to \(x=\log|f|=-u\) reverses every displayed
tropical coordinate.  To exchange min and max one also negates the value of
the tropical polynomial: \(F_{\max}(x)=-F_{\min}(-x)\).

The first uses are \Cref{chap:valuations} for graded reduction,
\Cref{chap:tropical-spaces} for calibrated polyhedra,
\Cref{chap:analytic-tropicalization} for analytic multiplicities,
and \Cref{chap:tropical-spectra-complex} for tropical spectra.

\section{Idempotent semirings}
\label{sec:appG-idempotent}

\begin{definition}[Idempotent semiring]
\label{def:appG-idempotent}
A commutative semiring
\((S,\oplus,\odot,\mathbb0,\mathbb1)\) is idempotent if
\[
                              a\oplus a=a
\tag{G.1}\label{eq:appG-idempotence}
\]
for every \(a\in S\).
Its natural order is
\[
                       a\preceq b
                       \quad\Longleftrightarrow\quad
                       a\oplus b=b.
\tag{G.2}\label{eq:appG-natural-order}
\]
\end{definition}

Reflexivity is idempotence.  If \(a\oplus b=b\) and \(b\oplus a=a\),
commutativity gives \(a=b\).  If also \(b\oplus c=c\), then
\(a\oplus c=a\oplus(b\oplus c)=(a\oplus b)\oplus c=c\), proving
transitivity.  The same identities show that \(a\oplus b\) is the least
upper bound of \(a,b\); distributivity makes multiplication order preserving.

\begin{example}[Max-plus and min-plus]
\label{ex:appG-max-min}
The max-plus semifield is
\[
 \mathbf T_{\max}
 =(\mathbf R\cup\{-\infty\},\max,+,-\infty,0),
\tag{G.3}\label{eq:appG-max-plus}
\]
and the min-plus semifield is
\[
 \mathbf T_{\min}
 =(\mathbf R\cup\{+\infty\},\min,+,+\infty,0).
\tag{G.4}\label{eq:appG-min-plus}
\]
Negation identifies the two.
The natural order of \(\mathbf T_{\max}\) is the usual order, while that of
\(\mathbf T_{\min}\) is the reverse of the usual order.
\end{example}

\begin{warning}{Do not adjoin both infinities}
\label{warn:appG-two-infinities}
The expression \(+\infty+(-\infty)\) is undefined.
Thus \(\mathbf R\cup\{-\infty,+\infty\}\), with tropical multiplication
given by ordinary addition, is not a semifield without an extra convention.
We use one-sided tropical extensions and handle the other endpoint
separately.
\end{warning}

\begin{definition}[Semimodule]
\label{def:appG-semimodule}
An \(S\)-semimodule is a commutative idempotent monoid
\((M,\oplus,\mathbb0)\) with an \(S\)-action satisfying the usual
distributive and unit axioms.
A tropical linear combination is a finite sum
\[
                         \bigoplus_i a_i\odot x_i.
\tag{G.5}\label{eq:appG-linear-combination}
\]
\end{definition}

Over \(\mathbf T_{\max}\), \eqref{eq:appG-linear-combination} is the
coordinatewise maximum of translates \(a_i+x_i\).
There is generally no additive inverse, so kernels and quotients are
described by congruences rather than by subtracting submodules.

\section{Valuations and tropical characters}
\label{sec:appG-valuations}

\begin{definition}[Nonnegative semivaluation]
\label{def:appG-semivaluation}
Let \(A\) be a ring.
A nonnegative semivaluation is a map
\[
                      v:A\longrightarrow\mathbf R_{\geq0}\cup\{+\infty\}
\]
such that
\[
\begin{aligned}
 v(0)&=+\infty,&v(1)&=0,\\
 v(fg)&=v(f)+v(g),&
 v(f+g)&\geq\min\{v(f),v(g)\}.
\end{aligned}
\tag{G.6}\label{eq:appG-semivaluation}
\]
\end{definition}

Its support
\[
                         \mathfrak p_v=\{f:v(f)=+\infty\}
\tag{G.7}\label{eq:appG-support}
\]
is prime.
Indeed, it is an ideal by the ultrametric inequality and nonnegativity, and
\(fg\in\mathfrak p_v\) forces one of \(v(f),v(g)\) to be infinite.
Since the support is prime, the quotient \(A/\mathfrak p_v\) is a domain.
The semivaluation extends to its fraction field by
\(v(\bar f/\bar g)=v(f)-v(g)\) for \(\bar g\neq0\).

\begin{definition}[Ideal semiring]
\label{def:appG-ideal-semiring}
Let \(\mathscr I_{\mathrm{fg}}(A)\) be the finitely generated ideals of
\(A\), with
\[
                         I\oplus J=I+J,\qquad
                         I\odot J=IJ.
\tag{G.8}\label{eq:appG-ideal-semiring}
\]
Its additive zero is \((0)\), its multiplicative unit is \(A\), and its
addition is idempotent.
\end{definition}

\begin{theorem}\label{thm:appG-ideal-characters}
The assignments
\[
\begin{aligned}
 v&\longmapsto\chi_v,\qquad
 \chi_v((f_1,\ldots,f_r))=\min_i v(f_i),\\
 \chi&\longmapsto v_\chi,\qquad
 v_\chi(f)=\chi((f))
\end{aligned}
\tag{G.9}\label{eq:appG-character-correspondence}
\]
give inverse bijections between nonnegative semivaluations on \(A\) and
unital semiring homomorphisms
\[
                 \chi:\mathscr I_{\mathrm{fg}}(A)\longrightarrow
                 \mathbf T_{\min}.
\]
With pointwise topologies on both sides, the bijection is a homeomorphism.
\end{theorem}

\begin{proof}
If \(g=\sum_ia_if_i\) and \(v(a_i)\geq0\), then
\[
                           v(g)\geq\min_iv(f_i).
\]
Thus the minimum in \eqref{eq:appG-character-correspondence} is independent
of the chosen generators.
Generators \(f_i g_j\) of \(IJ\) give
\[
 \chi_v(IJ)=\min_{i,j}(v(f_i)+v(g_j))
           =\chi_v(I)+\chi_v(J),
\]
while the union of generating sets gives
\(\chi_v(I+J)=\min\{\chi_v(I),\chi_v(J)\}\).

Conversely,
\((fg)=(f)(g)\) gives multiplicativity.
The identity \((f)+A=A\) yields
\[
                    0=\min\{v_\chi(f),0\},
\]
so \(v_\chi(f)\geq0\).
Since \((f+g)\subseteq(f,g)\), the natural semiring order gives
\[
 v_\chi(f+g)\geq\chi((f,g))
              =\min\{v_\chi(f),v_\chi(g)\}.
\]
The two constructions are inverse because principal ideals generate all
values in the first formula.
Finally, the inverse uses principal-ideal evaluations, while each forward
coordinate is a finite minimum of evaluations.  Finite minimum is
continuous on the one-sided extended real line, proving the topological claim.
\end{proof}

\begin{remark}[Noetherian and coherent versions]
\label{rem:appG-coherent-version}
If \(A\) is Noetherian, every ideal is finitely generated.
For a complex space, coherent ideal sheaves replace finitely generated
ideals.
The global tropical spectrum of
\Cref{chap:tropical-spectra-complex} uses exactly this coherent
version; it is not obtained by pretending a nonalgebraic complex space has
one global coordinate ring.
\end{remark}

\begin{definition}[Tropical spectrum]
\label{def:appG-tropical-spectrum}
For an idempotent semiring \(S\), its min-plus tropical spectrum is
\[
 \operatorname{tSpec}(S)
 :=
 \operatorname{Hom}_{\mathrm{Semiring}}(S,\mathbf T_{\min}),
\tag{G.10}\label{eq:appG-tropical-spectrum}
\]
equipped with the subspace topology from
\(\mathbf T_{\min}^{S}\).
\end{definition}

Compactness of a particular tropical spectrum is not formal from this
definition.
It requires closedness of the homomorphism equations and compact bounds on
all evaluation coordinates.
\Cref{chap:tropical-spectra-complex} proves those facts for coherent
ideal semirings on compact complex spaces.

\subsection{Extended tropical affine space}
\label{sec:appG-extended-affine}

Give \(\mathbf R\cup\{+\infty\}\) the topology transported from
\([0,+\infty)\) by \(u\mapsto \rme^{-u}\), with \(\rme^{-\infty}=0\).
The extended tropical affine space is
\[
                  \overline{\mathbf R}^{\,r}_{\min}
                  =(\mathbf R\cup\{+\infty\})^r.
\tag{G.11}\label{eq:appG-extended-affine}
\]
Its stratum indexed by \(J\subseteq\{1,\ldots,r\}\) has
\(u_j=+\infty\) exactly for \(j\in J\).
These are support strata: \(u_j=+\infty\) means that the corresponding
coordinate has seminorm zero.

An integer-affine formula extends across a support stratum only when all
its extended-real sums are defined there.
In particular, a Laurent expression containing both positive and negative
multiples of coordinates equal to \(+\infty\) cannot be evaluated by
formal cancellation.
One must restrict to a compatible toric stratum or change charts.

\section{Multiplicative seminorms and graded reduction}
\label{sec:appG-graded-reduction}

Let \(A\) carry a multiplicative non-Archimedean seminorm
\(|\cdot|\).
For \(r>0\), put
\[
 A_{\leq r}=\{f:|f|\leq r\},\qquad
 A_{<r}=\{f:|f|<r\}.
\tag{G.12}\label{eq:appG-norm-filtration}
\]

\begin{definition}[Graded reduction]
\label{def:appG-graded-reduction}
The graded reduction is
\[
 \widetilde A^\bullet
 =
 \bigoplus_{r>0}A_{\leq r}/A_{<r}.
\tag{G.13}\label{eq:appG-graded-reduction}
\]
For \(f\) with \(|f|=r>0\), its initial form
\(\widetilde f\) is its class in degree \(r\).
\end{definition}

\begin{proposition}\label{prop:appG-initial-calculus}
If \(|f|,|g|>0\), then
\[
                    \widetilde{fg}=\widetilde f\,\widetilde g.
\tag{G.14}\label{eq:appG-initial-product}
\]
If \(|f|=|g|=r\), write \([h]_r\) for the degree-\(r\) class of an
element \(h\in A_{\leq r}\).  Then
\[
 [f+g]_r=\widetilde f+\widetilde g=
 \begin{cases}
 \widetilde{f+g},&|f+g|=r,\\
 0\text{ in degree }r,&|f+g|<r.
 \end{cases}
\tag{G.15}\label{eq:appG-initial-sum}
\]
\end{proposition}

\begin{proof}
Multiplicativity gives \(|fg|=|f||g|\), and changing a representative by a
strictly smaller term changes the product by a strictly smaller term.
For sums, the ultrametric inequality puts \(f+g\) in \(A_{\leq r}\).
Its degree-\(r\) class is the sum of the two classes, which vanishes exactly
when cancellation lowers the norm.
\end{proof}

\begin{example}[Cancellation and the new degree]
\label{ex:appG-cancelled-degree}
If the norm drops but \(f+g\neq0\), its initial form in its new degree
need not vanish.  For example, on \(\mathbf Q((T))\) with
\(|T|=\rme^{-1}\), take \(f=1\) and \(g=-1+T\).  Their degree-one
classes cancel, while \(f+g=T\) has a nonzero initial form in degree
\(\rme^{-1}\).
\end{example}

In additive notation, for a semivaluation \(v\) with support
\(\mathfrak p_v\), set
\[
\begin{aligned}
 \mathcal P_\gamma&=\{f:v(f)\geq\gamma\},\\
 \mathcal P_\gamma^+&=\{f:v(f)>\gamma\}.
\end{aligned}
\]
Then
\[
 \operatorname{gr}_v(A/\mathfrak p_v)
 =
 \bigoplus_{\gamma}
 \mathcal P_\gamma/\mathcal P_\gamma^+.
\tag{G.16}\label{eq:appG-additive-graded}
\]
The changes of variables \(r=\rme^{-\gamma}\) and \(v=-\log|\cdot|\) identify
\eqref{eq:appG-graded-reduction} and
\eqref{eq:appG-additive-graded}.

\begin{warning}{Ordinary residue loses radii}
\label{warn:appG-residue-radii}
The degree-one part alone records only elements of norm at most one modulo
elements of norm less than one.
Different radii, and hence rational ranks used in tropical dimension
arguments, disappear.
The analytic multiplicities in
\Cref{chap:analytic-tropicalization} require the full graded
reduction, including residue-field degrees.
\end{warning}

For an analytic point \(x\), the relevant object is the graded residue field
of the completed residue field \(\mathcal H(x)\).
Its functorial and dimension-theoretic properties are due to Temkin; see
\cite{Temkin2004}.
This appendix uses only the elementary filtration calculus above.

\section{Initial forms and tropical hypersurfaces}
\label{sec:appG-initial-tropical}

Let \(K\) be a valued field, with \(v=-\log|\cdot|\), let \(M\) be a
lattice and \(N=M^\vee\), and consider a nonzero Laurent polynomial over \(K\)
\[
                              f=\sum_{m\in S}a_m\chi^m
\tag{G.17}\label{eq:appG-laurent}
\]
whose finite support \(S\subset M\) has been collected by exponent, with
\(a_m\neq0\) for every \(m\in S\).
For \(u\in N_{\mathbf R}\), define
\[
 \operatorname{trop}(f)(u)
 =
 \min_{m\in S}\{v(a_m)+\langle m,u\rangle\}.
\tag{G.18}\label{eq:appG-tropical-polynomial}
\]
It is a concave piecewise affine function in the min convention.

\begin{definition}[Corner locus]
\label{def:appG-corner-locus}
The tropical hypersurface of the expression \(f\) is the locus where the
minimum in \eqref{eq:appG-tropical-polynomial} is achieved by at least two
exponents.
\end{definition}

\begin{proposition}\label{prop:appG-initial-form-test}
At a weight \(u\), the weighted initial form
\(\operatorname{in}_u(f)\) is the sum, in the appropriate graded residue
ring, of the terms achieving the minimum in
\eqref{eq:appG-tropical-polynomial}.
The point \(u\) lies in the corner locus if and only if
\(\operatorname{in}_u(f)\) is not a monomial.
\end{proposition}

\begin{proof}
Give the coefficient initial form of \(a_m\) its valuation degree
\(v(a_m)\), and give \(\chi^m\) the additional degree
\(\langle m,u\rangle\) in the weighted graded Laurent algebra.
At the minimum total degree, retain precisely those nonzero coefficient
initial forms whose total degrees attain the minimum.
This construction does not require an element of \(K\) realizing the
real weight \(\langle m,u\rangle\).
Because the polynomial was collected by exponent, one surviving exponent
gives a monomial and at least two give a nonmonomial initial form.
\end{proof}

\begin{warning}{Fundamental-theorem boundary}
\label{warn:appG-fundamental-boundary}
\Cref{prop:appG-initial-form-test} is algebraic bookkeeping.
Identifying the corner locus with the closure of valuations of actual zeros
is the fundamental theorem of tropical geometry and requires hypotheses on
the valued field and the chosen closure.
It is not proved by the initial-form calculation alone.
\end{warning}

\section{Integer-affine and polyhedral conventions}
\label{sec:appG-polyhedra}

Let \(\Gamma\subset\mathbf R\) be an additive subgroup containing the value
group relevant to the chart.

\begin{definition}[Affine conventions]
\label{def:appG-affine-conventions}
\begin{enumerate}
\item An \emph{integer-affine function}, in the sense of
      Mikhalkin--Rau, is
      \(u\mapsto\langle m,u\rangle+c\) with \(m\in M\) and
      \(c\in\mathbf R\).  It is \(\Gamma\)-integer-affine when
      \(c\in\Gamma\).
\item A \emph{rational polyhedron} is an intersection of finitely many
      half-spaces
      \[
                         \langle m_i,u\rangle\geq c_i,
                         \qquad m_i\in M,\ c_i\in\mathbf R.
\tag{G.19}\label{eq:appG-polyhedron}
      \]
      It is \(\Gamma\)-rational when every \(c_i\) lies in \(\Gamma\).
\item An integer-affine map has lattice-linear part and arbitrary real
      translation; the prefix \(\Gamma\)- restricts the translation to the
      stated constant group.
\end{enumerate}
\end{definition}

Thus the unqualified terminology in this appendix is exactly that of
\cite[\S\S1.3 and 2.1]{MR18Tropical}; the qualifier \(\Gamma\) records the
additional arithmetic condition coming from a valued field.

Clearing denominators permits rational slopes, but it multiplies the raw
normal vector.  That multiple is not declared primitive: lattice indices
and calibration weights are always computed after returning to saturated
integral tangent lattices.

For a polyhedron \(\sigma\), let \(N_\sigma\) be the saturated lattice in
its tangent space.
If \(\tau\) is a codimension-one face of \(\sigma\), the primitive vector
\[
                         u_{\sigma/\tau}\in N_\sigma/N_\tau
\tag{G.20}\label{eq:appG-primitive-normal}
\]
points from \(\tau\) into \(\sigma\).

\begin{definition}[Weighted balancing]
\label{def:appG-balancing}
A pure \(q\)-dimensional weighted complex is balanced when for every
codimension-one cell \(\tau\),
\[
                  \sum_{\sigma\supset\tau}
                  m_\sigma u_{\sigma/\tau}=0
                  \quad\text{in }N/N_\tau.
\tag{G.21}\label{eq:appG-balancing}
\]
A tropical \(q\)-cycle is an equivalence class, under common weighted
refinement, of balanced pure \(q\)-dimensional rational polyhedral
complexes with nonzero integral weights.  A representative with positive
weights defines a tropical subspace.
\end{definition}

\begin{example}[Tropical line]
\label{ex:appG-tropical-line}
The rays generated by
\((1,0)\), \((0,1)\), and \((-1,-1)\), each of weight one, satisfy
\eqref{eq:appG-balancing}.
Replacing the last primitive vector by \((-2,-2)\) without changing its
weight fails the primitive-vector convention.
\end{example}

\section{Calibrations and signs}
\label{sec:appG-calibrations}

\begin{definition}[Calibration]
\label{def:appG-calibration}
An orientation of a \(q\)-cell \(\sigma\) is a generator of
\(\bigwedge^qN_\sigma\) up to positive sign.
A calibration assigns the oriented multivector
\[
                              \mu_\sigma=m_\sigma e_\sigma
\tag{G.22}\label{eq:appG-calibration}
\]
to every top cell, compatibly with subdivision.
\end{definition}

The calibration contains three distinct data:
the saturated tangent lattice, an orientation, and a multiplicity.
For an integer-affine map \(F\) that has rank \(q\) on \(\sigma\), the
pushforward weight on its image contains the lattice index
\[
                  [N_{F(\sigma)}:F(N_\sigma)].
\tag{G.23}\label{eq:appG-lattice-index}
\]
Summing these indexed weights over cells mapping to the same image cell is
the calibrated pushforward.

\begin{proposition}\label{prop:appG-calibration-tests}
\begin{enumerate}
\item Subdivision preserves a calibrated integral: new top cells inherit the
      old calibration, and new internal codimension-one faces cancel in
      pairs.
\item Under a generically finite integer-affine map \(F\) with locally
      finite pushed-forward calibration and properness on the support
      of the test form, calibrated
      pushforward satisfies \(\int F^*\alpha=\int\alpha\) for a
      compactly supported target \((q,q)\)-superform \(\alpha\), with
      the right integral taken against the pushed-forward calibration.
      Each source weight is multiplied once by
      \eqref{eq:appG-lattice-index}.
\item Reversing orientation changes a top real form integral by a sign, but
      does not change the positive measure of a positive \((q,q)\)-form,
      because the primed and double-primed orientations reverse together.
\end{enumerate}
\par
\end{proposition}

\begin{proof}
On an internal codimension-one face, the primitive normal vectors of the two
adjacent new cells are opposite.
For projection, choose lattice bases and let \(D\) be the absolute
determinant of the full-rank tangent lattice map.  Pullback of a
\((q,q)\)-superform contributes \(D^2\), one determinant from each
exterior factor.  Changing variables in the ordinary integral contributes
\(D^{-1}\), leaving exactly the index \(D\) in
\eqref{eq:appG-lattice-index}.  Sum this calculation over the locally
finite source cells.
The last assertion follows because the two exterior factors contribute the
same orientation sign, whose square is one.
\end{proof}

\begin{warning}{Inward versus outward}
\label{warn:appG-boundary-orientation}
Our vector \(u_{\sigma/\tau}\) points into \(\sigma\), as in
\eqref{eq:appG-primitive-normal}.
An outward-normal-first Stokes convention uses its negative.
Changing that boundary sign does not change the balancing equation
\eqref{eq:appG-balancing}; it changes the translation between balancing and
the boundary operator.  For the two labelled exterior factors, the exact
boundary traces are \eqref{eq:superform-boundary-traces}; in particular the
labelled contraction is determined by bidegree, not by a choice of words
such as \emph{first} or \emph{second} factor.
\end{warning}

\begin{example}[Lattice-index test]
\label{ex:appG-index-test}
In \(\mathbf R^2\), the primitive directions \((1,0)\) and \((1,2)\)
span a sublattice of index
\[
                         |\det((1,0),(1,2))|=2.
\]
Their transverse stable intersection has multiplicity two.
A calculation giving one has discarded the lattice calibration.
\end{example}

\section{Analytic tropical multiplicities}
\label{sec:appG-analytic-multiplicity}

Work over a complete non-Archimedean ground field.
Let \(Z\) be a pure \(q\)-dimensional effective analytic cycle on a good Hausdorff
boundaryless analytic space, and let
\(f:U\to\mathbf G_m^r\) be a moment map on a neighborhood of its support.
Assume that \(f\) has tropical dimension \(q\) along \(Z\).  Work on the
convenient top cells where an auxiliary projection to a \(q\)-torus
has the finite flat local degree specified in \cite[\S14.3.2]{CLD25}.
Denote by \(\Sigma_f(Z)\) the
\(q\)-dimensional characteristic skeletal locus of the moment map on \(Z\),
and choose a convenient cellular decomposition of it in the sense of
\cite[\S14.3.2]{CLD25}.  Let \(C\) be a top cell and put
\(\sigma=f_{\mathrm{trop}}(C)\).
Choose a torus homomorphism \(p:\mathbf G_m^r\to\mathbf G_m^q\) whose
tropicalization is injective on the affine span of \(\sigma\), let
\(s_p:\mathbf R^q\to\langle\sigma\rangle\) be the resulting affine inverse,
and let \(e_q\) be the lattice volume vector on \(\mathbf R^q\).  On the
relative interior of \(C\), the analytic local degree
\(d_C=\deg_x(p\circ f|_Z)\) is a constant positive integer, with degree extended additively
over the effective cycle \(Z\).  The canonical calibration is
\[
\begin{aligned}
 \mu_C&=d_C\,(s_p)_*e_q,\\
 \text{integer-affine pushforward factor}
   &=[N_{F(\sigma)}:F(N_\sigma)].
\end{aligned}
\tag{G.24}\label{eq:appG-analytic-multiplicity}
\]
The first line is independent of \(p\).  The local analytic degree, rather
than naive cardinality of a fibre, is essential over a nonalgebraically
closed field or in a non-strict analytic space.  Formula
\eqref{eq:appG-analytic-multiplicity} makes no additional identification with
an initial-degeneration multiplicity.  If several source cells map onto the
same target cell, their pushed-forward calibrations are added.  After choosing
the lattice orientation \(e_\sigma\), write the resulting target calibration
as \(\mu_\sigma=m(\sigma)e_\sigma\); these are the positive calibration
weights used below (and need not be integral without additional strictness and
ground-field hypotheses).  This target construction is used on
polyhedral subcharts where the map of calibrated skeletal supports is
topologically compact, so that the sums of source calibrations are locally
finite.  A global target statement additionally assumes this compactness
and properness of \(f\) on the support of \(Z\).

\begin{theorem}\label{thm:appG-analytic-calibration}
In the preceding setting, the target calibration is balanced on such
polyhedral subcharts at faces whose preimages lie in the boundaryless
interior.  Under the global compactness and properness conditions just
stated, the locally finite \(q\)-dimensional tropicalization
\(\operatorname{Trop}_f(Z)\) is balanced with the multiplicities
\eqref{eq:appG-analytic-multiplicity}.
These calibrations are compatible with product refinements and with
generically finite integer-affine projections on polyhedral subcharts for
which the induced map on the calibrated support is topologically compact; a
global pushforward additionally requires properness on the support.
\end{theorem}

\begin{proof}
We first state the imported foundational input exactly.  For an analytic
space of level \(q\), the local-degree construction and its independence of
the auxiliary torus projection are
\cite[Proposition~14.3.1 and \S14.3.2]{CLD25}; gluing and uniqueness are
\cite[Lemma~14.3.4 and Definition~14.3.5]{CLD25}.  For an effective cycle,
apply the construction to each pure component and add with its cycle
coefficient.  This gives precisely the first line of
\eqref{eq:appG-analytic-multiplicity}.  We use this cited construction as an
input and do not replace it by an unproved initial-degeneration formula.

The imported input for balancing is
\cite[Theorem~15.6.1]{CLD25}: the interior of an analytic space of level
\(q\), with the canonical calibration just described, is harmonious.  By the
definition of harmony, the calibrated discordance vector at every
codimension-one face \(\tau\) is zero.  Pair that vector with an integer
covector \(\ell\) vanishing on \(N_\tau\).  Using the primitive normal
\(u_{\sigma/\tau}\) pointing into each incident top cell
\(\sigma\supset\tau\), the resulting scalar identity is
\[
       \sum_{\sigma\supset\tau}
       m(\sigma)\,\ell(u_{\sigma/\tau})=0.
\]
Since integer covectors span
\((N/N_\tau)^\vee_{\mathbf R}\), the displayed scalar identities are
equivalent to the vector balancing equation
\eqref{eq:appG-balancing}.

Finally replace \(f\) by a product moment map defining a common refinement.
Both calibrations restrict to the canonical calibration on every common
top-dimensional subcell, by the uniqueness in
\cite[Lemma~14.3.4]{CLD25}; summing the refined integrals therefore recovers
the old integral.  For a generically finite integer-affine projection whose
map on the calibrated support is topologically compact, pushforward sums
the pushed-forward cell volume vectors.  When the projection is induced
by a homomorphism of moment tori, choose an auxiliary projection to
\(\mathbf G_m^q\) on the target and its composite on the source.
Both calibration computations then use the same finite flat moment map;
their affine inverses commute with the projection.  The local formula of
\cite[\S14.3.2]{CLD25} therefore gives compatibility.  An affine translation
only changes chart coordinates.  The related formula for finite analytic
morphisms between spaces of level \(q\),
\cite[Proposition~14.4.2]{CLD25}, uses the calibrated degree; this must not
be confused with the integer finite flat degree \(d_C\) above.
Expressing lattice-normalized measure in source and target bases contributes exactly
\([N_{F(\sigma)}:F(N_\sigma)]\), by
\Cref{prop:appG-calibration-tests}.  This proves the two asserted
compatibilities and also explains the properness clause in the statement.
\end{proof}

\begin{example}[Degree and inverse lattice index]
\label{ex:appG-degree-index}
For the identity moment on an analytic one-dimensional torus, use the
auxiliary projection \(p(t)=t^2\).  Its finite flat degree is two, even
in residue characteristic two.  The tropical map is \(u\mapsto2u\),
so its inverse sends the lattice vector to one half of that vector.
Thus \(d_C(s_p)_*e_1=2\cdot\tfrac12e_1=e_1\), agreeing with the
identity projection.  The integer local degree and the coefficient of
the calibration are distinct data.
\end{example}

\begin{example}[An infinite target fibre]
\label{ex:appG-infinite-target-fibre}
Take a countable disjoint union of analytic one-dimensional tori and let
\(f\) be the identity moment on every component.  The source is good,
Hausdorff and boundaryless, and each local calibration has weight one.
Every target skeletal point has infinitely many source preimages, however,
so summing the target weights gives infinity.  The compactness condition
in \Cref{thm:appG-analytic-calibration} excludes this example; local
harmony alone does not make a global pushed-forward calibration finite.
\end{example}

The goodness and boundarylessness hypotheses are part of the statement.
For spaces with boundary or without the goodness assumptions, the chapter
uses the corresponding local or compact-support formulation rather than
silently applying this global version.

\section{Cohomology retained by tropicalization}
\label{sec:appG-cohomologically-tropical}
\index{tropical cohomology}\index{cohomologically tropical variety}

The calibrated support records multiplicities and lattice directions.
Recovering the cohomology of an algebraic variety requires information
about its boundary strata as well. A precise comparison is available for
complex subvarieties of tori \cite{AAPS25Cohomologically}.
Here the complex ground field is trivially valued, and every cohomology
group has rational coefficients.

Let $V\subset T=(\mathbf C^*)^r$ be a closed subvariety and put
$S=\operatorname{Trop}(V)$. Choose a unimodular fan $\Sigma$ with support
$S$. Let $P_\Sigma$ be its complex toric variety, let $\overline V$ be
the closure of $V$ in $P_\Sigma$, and let $\overline S$ be the canonical
tropical compactification in the corresponding tropical toric variety.
For a cone $\sigma\in\Sigma$, write $V_\sigma$ for the intersection
with the torus orbit, and $S_\sigma$ for its tropicalization, the support
of the star fan. Bars denote their closures in the two compactifications.

The comparison of \cite[Definition~3.1 and \S2.7]{AAPS25Cohomologically}
is the graded ring homomorphism
\[
 \tau_\Sigma^*:H^\bullet_{\mathrm{trop}}(\overline S)
          \longrightarrow H^\bullet(\overline V;\mathbf Q),
 \qquad
 H^m_{\mathrm{trop}}=\bigoplus_{p+q=m}H^{p,q}_{\mathrm{trop}}.
\]
On $H^{k,k}_{\mathrm{trop}}(\overline S)$, it uses the identification
with the rational Chow group $A^k(P_\Sigma)_{\mathbf Q}$, followed by
the cycle-class map and restriction to $\overline V$. It sends
off-diagonal classes to zero. The construction also applies to each star
fan and its algebraic stratum. The variety is \emph{cohomologically
tropical} with respect to $\Sigma$ when
\[
 H^\bullet_{\mathrm{trop}}(\overline{S_\sigma})
       \xrightarrow{\ \tau_\sigma^*\ }
 H^\bullet(\overline{V_\sigma};\mathbf Q)
\]
is an isomorphism for every $\sigma\in\Sigma$, including the zero cone
\cite[Definition~1.1]{AAPS25Cohomologically}.

The term \emph{sch\"on} means that all the open algebraic strata are
nonsingular. The stronger term \emph{wundersch\"on} additionally requires
that every such stratum is connected and that $H^k(V_\sigma;\mathbf Q)$
is pure of weight $2k$ for every $k$
\cite[Definition~1.2 and \S2.2]{AAPS25Cohomologically}.
A $d$-dimensional tropical space is a \emph{tropical homology manifold}
when cap product with its fundamental class gives
\[
 H^{p,q}_{\mathrm{trop}}(U)
       \simeq H^{\mathrm{BM}}_{d-p,d-q}(U)
\]
for every open subset $U$, using tropical Borel--Moore homology on the
right \cite[Definition~2.6]{AAPS25Cohomologically}.

For a sch\"on $V$, \cite[Theorem~6.1]{AAPS25Cohomologically} states
that the comparison holds on all closed strata exactly when $V$ is
wundersch\"on and $S$ is a tropical homology manifold. Theorems~4.4
and~4.5 of the same source make these two algebraic properties independent
of the chosen unimodular fan. Its Theorem~6.2 also gives
\[
 H^k_{\mathrm{trop}}(S_\sigma)
       \simeq H^k(V_\sigma;\mathbf Q)
\]
on the open strata. The proof compares the algebraic Deligne weight
sequences with the tropical Deligne resolutions, using compatibility of
Gysin maps; the compactified comparison and local duality supply the
needed injectivity. These are the published comparison results, with the
sch\"on and all-strata hypotheses retained.

\subsection{Why the hypotheses involve every stratum}

The examples in \cite[\S8.1]{AAPS25Cohomologically} distinguish the
separate conditions. For generic complex coefficients, the curve
\[
                 a+bx+cy+dxy=0\quad\text{in }(\mathbf C^*)^2
\]
has compactification $\mathbf P^1$ and four punctures. It is
wundersch\"on, but its tropicalization is the four-ray coordinate cross.
The independent balancing relations on the two axes prevent local
tropical duality. In the canonical compactification, the tropical
cohomology in total degree two has dimension two, whereas
$H^2(\mathbf P^1;\mathbf Q)$ has dimension one
\cite[\S8.1.1]{AAPS25Cohomologically}.

A generic conic in $(\mathbf C^*)^2$ has six punctures and tropicalizes
to a three-ray tropical line of weight two. Local tropical duality holds,
and the comparison on the whole compactification is an isomorphism over
$\mathbf Q$. Each boundary orbit, however, meets the conic in two
distinct points: its algebraic $H^0$ is $\mathbf Q^2$, while the
corresponding tropical closed stratum has $H^0=\mathbf Q$. Thus
connectedness of the open strata, and comparison on every closed stratum,
are essential \cite[\S8.1.2]{AAPS25Cohomologically}.
The nodal curve $(x-a)(y-b)=0$, with $ab\ne0$, gives a complementary
boundary: it is cohomologically tropical but not sch\"on, and its
coordinate cross is not a tropical homology manifold
\cite[\S8.1.4]{AAPS25Cohomologically}. The sch\"on assumption is
therefore essential to the stated equivalence.

Hyperplane-arrangement complements in their arrangement torus embeddings
provide a positive class: for an essential projective arrangement, the
tropicalization is the associated Bergman fan, and the algebraic
strata and local tropical duality satisfy the comparison hypotheses
\cite[\S8.2]{AAPS25Cohomologically}.

\subsection{Hodge numbers in a semistable degeneration}

There is also a family comparison relevant to
\Cref{chap:hybrid-spaces}. Work in the setting of
\cite[\S7]{AAPS25Cohomologically}: a family
$\pi:\mathcal V\to D^*$ of smooth projective subvarieties, with fibre
$V_t$, over the punctured complex disk has a semistable extension
obtained inside a toric degeneration, with regular total space and reduced
simple-normal-crossing central fibre. The chosen unimodular triangulation
identifies its open central strata with subvarieties of the associated
tori. Let $Z$ be the extended tropicalization of the generic fibre over
$\mathbf C((t))$. Assume that $Z$ is a tropical homology manifold and
that every open central stratum, in this specified torus embedding, is
wundersch\"on. Theorem~7.1 of the source then gives
\[
 H^{p,q}_{\mathrm{trop}}(Z)
       \simeq\operatorname{Gr}^{W}_{2p}H^{p+q}_{\lim},
 \qquad
 h^{p,q}(V_t)=\dim_{\mathbf Q}H^{p,q}_{\mathrm{trop}}(Z).
\]
The odd weight-graded pieces of the limiting mixed Hodge structure
vanish. Here $H_{\lim}$ is the limiting mixed Hodge structure of this
family; the statement compares its weight-graded groups and the nearby
Hodge numbers. Its proof compares the algebraic and tropical Steenbrink
sequences through the stratum comparisons.

For the polyhedral and hybrid Monge--Amp\`ere questions in
\Cref{prob:frontier-polyhedral-regularity,prob:frontier-hybrid-ytd}, a
polarization, an operator domain, admissible lifts and a measure comparison
must still be supplied. The cohomological results above do not supply
envelope regularity, orthogonality or metric convergence. The cited
paper's \S8.4 discusses the tropical Monge--Amp\`ere direction as a
further research programme; its announced differential-calculus sequel
is not an input to this appendix.

\section{Global sign and normalization table}
\label{sec:appG-sign-table}

\begin{center}
\begin{tabular}{p{.28\textwidth}p{.27\textwidth}p{.34\textwidth}}
\toprule
datum & convention & consequence\\
\midrule
valuation & \(v(f)=-\log|f|\) & sums tropicalize with \(\min\)\\
tropical coordinate &
\(u_i=-\log|f_i|\) &
monomials become \(\langle m,u\rangle+v(a_m)\)\\
max-plus coordinate &
\(\log|f_i|=-u_i\) &
\(F_{\max}(x)=-F_{\min}(-x)\); coordinate reflection alone does not switch
min and max\\
toric support function &
\(h_P(u)=\sup_{m\in P}\langle m,u\rangle\) &
metric recession uses \(h_P(-u)\) in
\eqref{eq:appJ-toric-recession}\\
tropical-cycle top-cell weight &
nonzero integer, positive for an effective tropical subspace &
orientation is separate from the sign of a cycle weight\\
analytic calibration weight &
positive real coefficient obtained from integer local degree and inverse
lattice map &
integrality requires the additional strictness and ground-field hypotheses
stated near \eqref{eq:appG-analytic-multiplicity}\\
stable intersection &
lattice index times weights &
nonprimitive determinants are not discarded\\
\bottomrule
\end{tabular}
\end{center}

\section{Normalization checks}
\label{sec:appG-unit-tests}

\begin{enumerate}
\item In \(\mathbf T_{\min}\),
      \(3\oplus5=3\), \(3\odot5=8\),
      \(\mathbb0=+\infty\), and \(\mathbb1=0\).
\item For a semivaluation,
      \(\chi_v((f,g))=\min\{v(f),v(g)\}\) and
      \(\chi_v((f)(g))=v(f)+v(g)\).
\item If two equal-norm terms cancel, their degree-\(r\) initial forms sum
      to zero and the norm drops, as in
      \eqref{eq:appG-initial-sum}.
\item For \(f=1+T\) over the trivially valued field,
      \(\operatorname{trop}(f)(u)=\min\{0,u\}\), whose corner is \(u=0\).
\item The three primitive rays of
      \Cref{ex:appG-tropical-line} balance.
\item The determinant-two intersection of
      \Cref{ex:appG-index-test} must retain multiplicity two after
      every subdivision.
\item Replacing \(u=-\log|\chi|\) by \(x=\log|\chi|=-u\), and at the same
      time replacing \(F_{\min}\) by
      \(F_{\max}(x)=-F_{\min}(-x)\), exchanges the min and max
      descriptions.
\end{enumerate}

The first two computations are the definitions
\eqref{eq:appG-min-plus} and \eqref{eq:appG-character-correspondence}.
The third is \Cref{prop:appG-initial-calculus}; the fourth follows by
substitution in \eqref{eq:appG-tropical-polynomial}; the next two are
\Cref{ex:appG-tropical-line,ex:appG-index-test}; and the last is the sign
identity stated at the beginning of the appendix.  Thus these are checked
normalizations.

\section*{Source note}
\label{sec:appG-source-note}

Graded reductions in analytic geometry follow Temkin
\cite{Temkin2004}.
Polyhedral spaces, canonical calibrations, and analytic balancing follow
\cite{CLD25}.
The complex all-strata cohomology comparison and its semistable
Hodge-number application follow the published
\cite{AAPS25Cohomologically}, in the setting of
\Cref{sec:appG-cohomologically-tropical}.
The unqualified integer-affine, rational-polyhedral, weighting, and
balancing terminology follows
\cite[\S\S1.3, 2.1 and 4.2]{MR18Tropical}.
The ideal-semiring tropical spectrum for complex spaces follows
\cite{MP24}; \Cref{thm:appG-ideal-characters} isolates its elementary
affine algebra.
The global comparison with projective trivially valued analytification is
proved in \Cref{chap:tropical-spectra-complex}, not in the abstract
semiring formalism.
